\documentclass[10pt,ignorenonframetext,,aspectratio=149]{beamer}
\usefonttheme{serif} % use mainfont rather than sansfont for slide text
\setbeamertemplate{caption}[numbered]
\setbeamertemplate{caption label separator}{: }
\setbeamercolor{caption name}{fg=normal text.fg}
\usepackage{lmodern}
\usepackage{amssymb,amsmath}
\usepackage{ifxetex,ifluatex}
\usepackage{fixltx2e} % provides \textsubscript
\ifnum 0\ifxetex 1\fi\ifluatex 1\fi=0 % if pdftex
  \usepackage[T1]{fontenc}
  \usepackage[utf8]{inputenc}
\else % if luatex or xelatex
  \ifxetex
    \usepackage{mathspec}
  \else
    \usepackage{fontspec}
  \fi
  \defaultfontfeatures{Ligatures=TeX,Scale=MatchLowercase}
  \newcommand{\euro}{€}
    \setmainfont[]{Open Sans}
\fi
% use upquote if available, for straight quotes in verbatim environments
\IfFileExists{upquote.sty}{\usepackage{upquote}}{}
% use microtype if available
\IfFileExists{microtype.sty}{%
\usepackage{microtype}
\UseMicrotypeSet[protrusion]{basicmath} % disable protrusion for tt fonts
}{}
\usepackage{graphicx,grffile}
\makeatletter
\def\maxwidth{\ifdim\Gin@nat@width>\linewidth\linewidth\else\Gin@nat@width\fi}
\def\maxheight{\ifdim\Gin@nat@height>\textheight0.8\textheight\else\Gin@nat@height\fi}
\makeatother
% Scale images if necessary, so that they will not overflow the page
% margins by default, and it is still possible to overwrite the defaults
% using explicit options in \includegraphics[width, height, ...]{}
\setkeys{Gin}{width=\maxwidth,height=\maxheight,keepaspectratio}

% Comment these out if you don't want a slide with just the
% part/section/subsection/subsubsection title:
\AtBeginPart{
  \let\insertpartnumber\relax
  \let\partname\relax
  \frame{\partpage}
}
\AtBeginSection{
  \let\insertsectionnumber\relax
  \let\sectionname\relax
  \frame{\sectionpage}
}
\AtBeginSubsection{
  \let\insertsubsectionnumber\relax
  \let\subsectionname\relax
  \frame{\subsectionpage}
}

\setlength{\emergencystretch}{3em}  % prevent overfull lines
\providecommand{\tightlist}{%
  \setlength{\itemsep}{0pt}\setlength{\parskip}{0pt}}
\setcounter{secnumdepth}{0}

\title{Some Least Squares Extensions}
\author{Henrique Veras}
\date{}

%% Here's everything I added.
%%--------------------------

\usepackage{graphicx}
\usepackage{rotating}
%\setbeamertemplate{caption}[numbered]
\usepackage{hyperref}
\usepackage{caption}
\usepackage[normalem]{ulem}
%\mode<presentation>
\usepackage{wasysym}
%\usepackage{amsmath}


% Get rid of navigation symbols.
%-------------------------------
\setbeamertemplate{navigation symbols}{}

% Optional institute tags and titlegraphic.
% Do feel free to change the titlegraphic if you don't want it as a Markdown field.
%----------------------------------------------------------------------------------
\institute{PIMES/UFPE}

% \titlegraphic{\includegraphics[width=0.3\paperwidth]{\string~/Dropbox/teaching/clemson-academic.png}} % <-- if you want to know what this looks like without it as a Markdown field. 
% -----------------------------------------------------------------------------------------------------


% Some additional title page adjustments.
%----------------------------------------
\setbeamertemplate{title page}[empty]
%\date{}
\setbeamerfont{subtitle}{size=\small}

\setbeamercovered{transparent}

% Some optional colors. Change or add as you see fit.
%---------------------------------------------------
\definecolor{clemsonpurple}{HTML}{522D80}
 \definecolor{clemsonorange}{HTML}{F66733}
\definecolor{uiucblue}{HTML}{003C7D}
\definecolor{uiucorange}{HTML}{F47F24}


% Some optional color adjustments to Beamer. Change as you see fit.
%------------------------------------------------------------------
\setbeamercolor{frametitle}{fg=clemsonpurple,bg=white}
\setbeamercolor{title}{fg=clemsonpurple,bg=white}
\setbeamercolor{local structure}{fg=clemsonpurple}
\setbeamercolor{section in toc}{fg=clemsonpurple,bg=white}
% \setbeamercolor{subsection in toc}{fg=clemsonorange,bg=white}
\setbeamercolor{footline}{fg=clemsonpurple!50, bg=white}
\setbeamercolor{block title}{fg=clemsonorange,bg=white}


\let\Tiny=\tiny


% Sections and subsections should not get their own damn slide.
%--------------------------------------------------------------
\AtBeginPart{}
\AtBeginSection{}
\AtBeginSubsection{}
\AtBeginSubsubsection{}

% Suppress some of Markdown's weird default vertical spacing.
%------------------------------------------------------------
\setlength{\emergencystretch}{0em}  % prevent overfull lines
\setlength{\parskip}{0pt}


% Allow for those simple two-tone footlines I like. 
% Edit the colors as you see fit.
%--------------------------------------------------
\defbeamertemplate*{footline}{my footline}{%
    \ifnum\insertpagenumber=1
    \hbox{%
        \begin{beamercolorbox}[wd=\paperwidth,ht=.8ex,dp=1ex,center]{}%
      % empty environment to raise height
        \end{beamercolorbox}%
    }%
    \vskip0pt%
    \else%
        \Tiny{%
            \hfill%
		\vspace*{1pt}%
            \insertframenumber/\inserttotalframenumber \hspace*{0.1cm}%
            \newline%
            \color{clemsonpurple}{\rule{\paperwidth}{0.4mm}}\newline%
            \color{clemsonorange}{\rule{\paperwidth}{.4mm}}%
        }%
    \fi%
}

% Various cosmetic things, though I must confess I forget what exactly these do and why I included them.
%-------------------------------------------------------------------------------------------------------
\setbeamercolor{structure}{fg=blue}
\setbeamercolor{local structure}{parent=structure}
\setbeamercolor{item projected}{parent=item,use=item,fg=clemsonpurple,bg=white}
\setbeamercolor{enumerate item}{parent=item}

% Adjust some item elements. More cosmetic things.
%-------------------------------------------------
\setbeamertemplate{itemize item}{\color{clemsonpurple}$\bullet$}
\setbeamertemplate{itemize subitem}{\color{clemsonpurple}\scriptsize{$\bullet$}}
\setbeamertemplate{itemize/enumerate body end}{\vspace{.6\baselineskip}} % So I'm less inclined to use \medskip and \bigskip in Markdown.

% Automatically center images
% ---------------------------
% Note: this is for ![](image.png) images
% Use "fig.align = "center" for R chunks

\usepackage{etoolbox}

\AtBeginDocument{%
  \letcs\oig{@orig\string\includegraphics}%
  \renewcommand<>\includegraphics[2][]{%
    \only#3{%
      {\centering\oig[{#1}]{#2}\par}%
    }%
  }%
}

% I think I've moved to xelatex now. Here's some stuff for that.
% --------------------------------------------------------------
% I could customize/generalize this more but the truth is it works for my circumstances.

\ifxetex
\setbeamerfont{title}{family=\fontspec{Titillium Web}}
\setbeamerfont{frametitle}{family=\fontspec{Titillium Web}}
\usepackage[font=small,skip=0pt]{caption}
 \else
 \fi

% Okay, and begin the actual document...

\begin{document}
\frame{\titlepage}

\hypertarget{econometrics}{%
\section{Econometrics}\label{econometrics}}

\hypertarget{intro}{%
\subsection{Intro}\label{intro}}

\begin{frame}{Intro}
In this lecture, we discuss some extensions of the least squares model.

\vfill

Specifically, we expand the analysis to incorporate

\begin{enumerate}
\tightlist
\item
  Qualitative information
\item
  Nonlinear relationships
\item
  Interactions between explanatory variables
\end{enumerate}
\end{frame}

\hypertarget{using-binary-variables}{%
\subsection{Using Binary Variables}\label{using-binary-variables}}

\begin{frame}{Binary Variables in Regression}
\protect\hypertarget{binary-variables-in-regression}{}
Consider the regression model

\[y_i=\mathbf{x}_i'\beta+\gamma d_i + \varepsilon\]

where \(d_i=1\) for some condition occurring and 0 if not. We call these
types of variables \textbf{dummy variables} or indicator variables.

\vfill

For example, consider the earnings regression once again, now
incorporating a qualitative information:

\[Earnings = \beta_1+\beta_2 Educ + \beta_3 Age + \beta_4 union + \varepsilon\]
where \(union=1\) if the worker is registered in a union and 0
otherwise.

\vfill

How do we interpret \(\beta_4\)?

\vfill
\end{frame}

\begin{frame}{Binary Variables in Regression}
\protect\hypertarget{binary-variables-in-regression-1}{}
Dummy variables are particularly useful in loglinear regressions.

\vfill

Consider the model

\[\ln y=\beta_1+\beta_2x +\beta_3 d + \varepsilon\]

\vfill

The coefficient \(d\) indicates a multiplicative shift of the function.

\vfill
\end{frame}

\begin{frame}{Binary Variables in Regression}
\protect\hypertarget{binary-variables-in-regression-2}{}
The percentage change in \(E[\mathbf{y}|\mathbf{X},d]\) associated with
the change in \(d\) is

\[\begin{split}
\%(\Delta E[\mathbf{y}|\mathbf{X},d]/\Delta d) & = 100\% \left(\frac{E[\mathbf{y}|\mathbf{X},d=1]-E[\mathbf{y}|\mathbf{X},d=0]}{E[\mathbf{y}|\mathbf{X},d=0]}\right) \\
& = 100\%[\exp(\beta_3)-1]
\end{split}\]

\vfill

Take the Monet paintings example again. We can introduce qualitative
information for a better model fit.

\vfill

On the data set, we also have information whether the painting was
signed by Monet or not (of the 430 sales in the sample, 77 are for
unsigned paintings). This definitely affects its price!

\vfill
\end{frame}

\begin{frame}{Monet's Paintings}
\protect\hypertarget{monets-paintings}{}
The table below shows the regression results from the model

\[\ln(Price)=\beta_1+\beta_2\ln(\text{Area})+\beta_3\text{Aspect Ratio}+\beta_4\text{Signature}+\varepsilon\]

\includegraphics{/Users/henriquefonseca/Desktop/temp/Rmarkdown-practice/henriqueveras.github.io/files/Econometrics/Lecture Notes/6/table_6.2.png}

\vfill
\end{frame}

\begin{frame}{Monet's Paintings}
\protect\hypertarget{monets-paintings-1}{}
Computing this result for a painting of the same area and aspect ratio,
we find the model predicts that the signature effect would be

\[100\% \times \frac{\Delta E(Price)}{Price}=100\%[\exp(1.26)-1]=256\%\]

The effect of a signature on an otherwise similar painting is to more
than double the price.

\vfill

Using the delta method, we obtain an estimated standard error for
\([\exp(b_3)^2 - 1]\) of the square root of
\([\exp(b_3)] \times 0.1253^2\) , which is 0.4417.

\vfill

For the percentage difference of 252\%, we have an estimated standard
error of 44.17\%.

\vfill
\end{frame}

\begin{frame}{Several Categories}
\protect\hypertarget{several-categories}{}
When there are several categories, a set of binary variables is
necessary.

\vfill

Consider a consumption function for quarterly data:

\[C_t=\beta_1 + \beta_2x_t + \delta_1 D_{1t}+\delta_2 D_{2t}+\delta_3 D_{3t}+\varepsilon_t\]

Notice that we only introduce 3 quarterly dummy variables, as the fourth
one would add one to all observations, adding to the intercept. This is
known as the \textbf{dummy variable trap}.

\vfill
\end{frame}

\begin{frame}{Several Categories}
\protect\hypertarget{several-categories-1}{}
The preceding is a means of \textbf{deseasonalizing} the data. Consider
the alternative formulation:

\[C_t=\beta x_t +  \delta_1 D_{1t}+\delta_2 D_{2t}+\delta_3 D_{3t}++  \delta_4 D_{4t}+\varepsilon_t\]

\vfill

The FWL theorem suggests that we can obtain \(\beta\) by regressing both
\(C\) and \(x\) on the four quarterly dummies and then regress the
residuals of the first on the residuals of the second.

\vfill

This means we are \emph{removing} the variation in both \(C\) and \(x\)
associated with each quarter of the year.

\vfill
\end{frame}

\begin{frame}{Threshold Effects and Categorical Variables}
\protect\hypertarget{threshold-effects-and-categorical-variables}{}
Suppose we are interested in estimating the effect of education on
earnings and the available data consists of the highest level of
educational attainment, such as less than high school (LTHS), high
school (HS), college (C), post-graduate (PG).

\vfill

A flexible model would use three (or four) binary variables, one for
each level of education:

\[Earnings = \beta_1+\beta_2 Age + \delta_{HS}HS+\delta_{C}C+\delta_{PG}PG+\varepsilon\]
How do we interpret these coefficients?

\vfill
\end{frame}

\begin{frame}{Threshold Effects and Categorical Variables}
\protect\hypertarget{threshold-effects-and-categorical-variables-1}{}
The correspondence between the coefficients and income for a given age
is

\[\begin{split}
\text{LTHS: }  E[Earnings|Age,LTHS] & = \beta_1 +\beta_2 Age \\
\text{HS: }  E[Earnings|Age,HS] & =  \beta_1 +\beta_2 Age+ \delta_{HS} \\
\text{C: }  E[Earnings|Age,C] & = \beta_1 +\beta_2 Age+ \delta_{C} \\
\text{PG: } E[Earnings|Age,PG] & = \beta_1 +\beta_2 Age+\delta_{PG} 
\end{split}\]

We can then estimate, for example, \(\delta_{PG}-\delta_{C}\), which
might be of interest.

\vfill
\end{frame}

\begin{frame}{Threshold Effects and Categorical Variables}
\protect\hypertarget{threshold-effects-and-categorical-variables-2}{}
Another way to estimate the relationship above is to recast the
variables in terms of whether the individual has some degree and not the
highest level attained.

\vfill

How would you redefine the above equations? How would we interpret each
dummy variable coefficient?

\vfill
\end{frame}

\begin{frame}{Transition Tables}
\protect\hypertarget{transition-tables}{}
Hodge and Shankar (2014) propose a useful framework for arranging the
effects of a sequence of categories:

\includegraphics{/Users/henriquefonseca/Desktop/temp/Rmarkdown-practice/henriqueveras.github.io/files/Econometrics/Lecture Notes/6/table_6.7.png}

\vfill

How would we calculate the standard error of the partial effects?
\end{frame}

\hypertarget{nonlinearities-in-the-variables}{%
\subsection{Nonlinearities in the
Variables}\label{nonlinearities-in-the-variables}}

\begin{frame}{Nonlinearities in the Variables}
Let \(\mathbf{z}=z_1,z_2,\cdots z_L\) be a set of \(L\) independent
variables; let \(f_1, f_2, \cdots f_K\) be \(K\) linearly independent
functions of \(\mathbf{z}\); let \(g(y)\) be an observable function of
\(y\);

\vfill

The linear regression model may be written

\[\begin{split}
g(y) & =  \beta_1f_1(\mathbf{z})+\beta_2f_2(\mathbf{z})+\cdots+\beta_Kf_K(\mathbf{z})+\varepsilon \\
 & = \beta_1x_1+\beta_2x_2+\cdots+\beta_Kx_K+\varepsilon \\
 & = \mathbf{x'\beta} + \varepsilon
\end{split}\]

This linear model can be tailored to any number of situations.

\vfill
\end{frame}

\begin{frame}{Functional Forms}
\protect\hypertarget{functional-forms}{}
A commonly used form of regression model is the loglinear model,

\[\ln y = \ln\alpha + \sum_k{\beta_k\ln x_k}+\varepsilon=\beta_1+ \sum_k{\beta_k\ln x_k}+\varepsilon\]

\vfill

In this model, the coefficients are elasticities:

\[\left(\frac{\partial y}{\partial X_k}\right)\left(\frac{X_k}{y}\right)=\frac{\partial \ln y}{\partial \ln X_k}=\beta_k\]

\vfill

A \textbf{semilog} model is

\[\ln y=\beta_1 +\beta_2 x+\varepsilon\]

\vfill

How do we interpret \(\beta_2\)?
\end{frame}

\begin{frame}{Functional Forms}
\protect\hypertarget{functional-forms-1}{}
An alternative approach sometimes taken is to measure the variables and
associated changes in standard deviation units:

\[x_{ik}^*=(x_{ik}-\bar{x}_k)/s_k \text{,   } y_{i}^*=(y_{i}-\bar{y})/s_y\]

\vfill

How do we interpret the \(\beta_2\) coefficients in a regression model
given by

\[y_i^*=\beta_1+\beta_2x_i^*+\varepsilon\]

\vfill
\end{frame}

\hypertarget{interaction-effects}{%
\subsection{Interaction Effects}\label{interaction-effects}}

\begin{frame}{Interaction Effects}
Another useful formulation of the regression model is one with
interaction terms.

\vfill

Suppose the following model:

\[\ln wage = \beta_1 Educ + \beta_2 female + \beta_3 Educ\times female+\cdots+\varepsilon\]

\vfill

In this model,

\[\frac{\partial E[\ln wage|Educ, fem, \cdots]}{\partial Educ}=\beta_1+\beta_3female\]

which implies that the marginal effect of education differs between men
and women.

\vfill
\end{frame}

\begin{frame}{Interaction Effects}
\protect\hypertarget{interaction-effects-1}{}
To perform hypothesis testing and constructing confidence intervals, the
necessary standard error is computed from the square root of

\[\begin{split}Var\left(\frac{\partial \hat{E[\ln wage|Educ, fem, \cdots]}}{\partial Educ}\right)= & Var[\hat{\beta}_1]+female^2Var[\hat{\beta}_3] \\
& +2female \text{ } Cov(\hat{\beta}_1,\hat{\beta}_3)\end{split}\]

\vfill
\end{frame}


\section[]{}
\frame{\small \frametitle{Table of Contents}
\tableofcontents}
\end{document}
